\section{Equilibrium Characterization}
\label{sec:equilibrium}

We solve the game by backward induction. Let
\begin{equation}
D\equiv4-\theta^2,
\qquad
m\equiv a-c>0,
\qquad
r_i\equiv x_i+\mu g_i+\nu h_i.
\end{equation}

\subsection{Stage 3: nonnegative Cournot continuation}

For an active firm, the quantity first-order condition is
\[
a-c_i-2q_i-\theta q_j=0.
\]
If both firms are active, solving the two equations gives
\begin{equation}
q_i
=\frac{(2-\theta)m+2r_i-\theta r_j}{D}.
\label{eq:cournot-output}
\end{equation}
On this branch $p_i-c_i=q_i$, so operating profit is $q_i^2$. This identity is not global: an inactive firm instead satisfies the corresponding Kuhn--Tucker inequality.

It is useful to write
\[
v_i\equiv m+x_i+\mu g_i+\nu h_i.
\]
The complete nonnegative Cournot continuation is then piecewise. Both firms are active when $2v_i>\theta v_j$ for each $i$. Firm $A$ is the active monopolist when $v_B\leq\theta v_A/2$, in which case $q_A=v_A/2$ and $q_B=0$; the symmetric condition applies for firm $B$. These boundary regimes matter because a government can otherwise profitably deviate far enough to change market structure. The maintained restriction $m>0$ rules out a separate both-inactive baseline region for the continuation map used below.

\subsection{Stage 2: investment and the active-duopoly branch}

On the two-active-firm branch, differentiating reduced operating profit $q_i^2$ gives
\begin{equation}
k_x x_i=\frac{4q_i}{D},
\qquad
k_g g_i=\frac{4\mu q_i}{D}+s_i.
\end{equation}
Hence
\begin{equation}
x_i^*=\frac{4q_i}{Dk_x},
\qquad
g_i^*=\frac{4\mu q_i/D+s_i}{k_g}.
\label{eq:investment-solutions}
\end{equation}
The investment Hessian on this branch is
\begin{equation}
H_F=
\frac{8}{D^2}
\begin{pmatrix}
1&\mu\\
\mu&\mu^2
\end{pmatrix}
-
\begin{pmatrix}k_x&0\\0&k_g\end{pmatrix}.
\label{eq:firm-hessian}
\end{equation}
Condition \eqref{eq:regularity} is sufficient for the required firm-stage concavity and reduced-system regularity.

Define
\begin{equation}
\lambda\equiv\frac{4R}{D},
\qquad
y_i\equiv\frac{\mu}{k_g}s_i+\nu h_i,
\qquad
w_i\equiv m+y_i,
\qquad
L\equiv2-\lambda.
\label{eq:lambda-y}
\end{equation}
The interior Stage-2 equilibrium quantities satisfy the compact system
\begin{equation}
Lq_i+\theta q_j=w_i.
\label{eq:reduced-output-system}
\end{equation}
Under \eqref{eq:regularity}, $L>\theta$ for all $\theta\in[0,1]$. Therefore
\begin{equation}
q_i=q_0+t_0y_i+t_1y_j,
\label{eq:q-linear-policies}
\end{equation}
where
\begin{align}
q_0&=\frac{m}{L+\theta},\\
t_0&=\frac{L}{L^2-\theta^2}>0,\\
t_1&=-\frac{\theta}{L^2-\theta^2}\leq0.
\label{eq:t0t1}
\end{align}
Thus either local policy raises the local firm's output through $y_i$, while the rival's policy lowers local output whenever $\theta>0$.

The coefficient
\begin{equation}
R=\underbrace{\frac1{k_x}}_{\text{conventional investment}}
+\underbrace{\frac{\mu^2}{k_g}}_{\text{green investment}}
\end{equation}
controls the strength of the investment feedback. Removing conventional competitive investment deletes the first term. For fixed primitives, $\lambda=4R/(4-\theta^2)$ rises with product substitutability.

\subsection{Full Stage-2 regime map}

To verify global equilibrium, the investment game must also be solved when one firm becomes inactive. Let $z_i=x_i+\mu g_i$ be private cost reduction and define $u_i=z_i-(\mu/k_g)s_i$. Conditional on $u_i$, the cost-minimizing investment composition is
\begin{equation}
x_i=\frac{u_i}{Rk_x},
\qquad
g_i=\frac{s_i}{k_g}+\frac{\mu u_i}{Rk_g}.
\label{eq:scalar-investment-composition}
\end{equation}
This reduces Stage 2 to a one-dimensional investment choice for each firm. To see why the reduction is valid beyond the duopoly branch, complete the square in $g_i-s_i/k_g$. If $u_i<0$, replacing the proposed investments by $(x_i,g_i)=(0,s_i/k_g)$ weakly raises the firm's product-market position while strictly reducing its investment cost; hence $u_i<0$ is strictly dominated. For $u_i\geq0$ and $\mu\geq0$, the unique minimum-cost nonnegative investment pair is given by \eqref{eq:scalar-investment-composition}. Up to the policy-dependent constant $s_i^2/(2k_g)$, the scalar firm objective is operating profit less $u_i^2/(2R)$. Its second derivative is $-1/R$ when the firm is inactive, $8/D^2-1/R$ on the duopoly branch, and $1/2-1/R$ in own monopoly. These are strictly negative under \eqref{eq:regularity}. At the firm's own entry boundary, the derivative of operating profit is zero on both sides; at the point where its investment renders the rival inactive, that derivative weakly falls from $4q_i/D$ to $q_i$. Thus the scalar problem is globally strictly concave across the connected pieces, not merely separately concave within each formula. Inactive production does not imply withdrawal from the investment/subsidy game: an inactive firm can still have $g_i=s_i/k_g$.

Under the maintained $m>0$ restriction and for $\theta>0$, the unique continuation is therefore classified by $w_A/w_B$:

\begin{itemize}
\item If $\theta w_A\geq(2-R)w_B$, firm $A$ is the active monopolist and
\[
q_A=\frac{w_A}{2-R},\qquad q_B=0,
\qquad u_A=Rq_A,\quad u_B=0.
\]
\item If $Lw_B\leq\theta w_A\leq(2-R)w_B$, firm $A$ limit-prices at the kink and
\[
q_A=\frac{w_B}{\theta},\qquad q_B=0,
\qquad u_A=\frac{2w_B}{\theta}-w_A,\quad u_B=0.
\]
\item If $\theta w_A<Lw_B$ and $\theta w_B<Lw_A$, both firms are active and \eqref{eq:reduced-output-system} applies, with $u_i=\lambda q_i$.
\item The remaining two regions are the symmetric $B$-kink and $B$-monopoly cases.
\end{itemize}

At the regime boundaries the continuation formulas coincide. At $\theta=0$ strategic exclusion is impossible and the two-active-firm continuation covers the relevant policy problem.

\subsection{Stage 1: government best responses and global qualification}

On the interior duopoly branch, substituting \eqref{eq:investment-solutions} and \eqref{eq:q-linear-policies} into \eqref{eq:government-welfare} makes each government's reduced objective quadratic in $(s_A,h_A,s_B,h_B)$. Let $z_i=(s_i,h_i)'$. Whenever the own-policy Hessian $H_i\equiv W_{z_i z_i}$ is negative definite, the branch-specific interior best response is unique and affine in the rival policy vector. Section~\ref{sec:results} characterizes its cross-instrument derivative exactly.

Branch-specific concavity alone is not sufficient for SPNE. For the canonical rational witness, the proof of the global comparison has four logically distinct parts. First, the symmetric policy candidate is nonnegative, supports strictly positive output for both firms, and has a negative-definite own-policy Hessian. Second, holding the rival at this policy, the exact inequality $Lm-\theta w_B>0$ holds throughout $\theta\in[0,1]$. Since an own nonnegative policy implies $w_A\geq m$, a unilateral deviation cannot reach a rival-kink or rival-monopoly continuation. Third, the remaining own-kink and own-monopoly branches meet at $\theta w_A=(2-R)w_B$. Maximizing welfare tangentially along this boundary yields a feasible positive $(s_A,h_A)$ and a zero tangent derivative. On the boundary the derivative toward increasing $w_A$ is positive from the kink side and negative from the monopoly side, while each branch is strictly concave; the common boundary candidate thus bounds both exclusionary branches. Finally, the symmetric duopoly candidate has strictly higher welfare than this best boundary challenger. The difference $\Delta W$ is positive for every $\theta\in(0,1]$, with representative gaps $7.649$ at $\theta=.5$, $.428$ at $.9$, and $.073$ at $1$.

Appendix~\ref{app:global-equilibrium} gives the explicit comparison map and handles the $\theta=0$ endpoint by continuously extending $\theta^2\Delta W(\theta)$. The boundary policies and one-sided normal derivatives are scaled by $\theta$ before taking that endpoint limit; at $\theta=0$ the product markets decouple and the exclusionary region is unavailable. Nonsingularity of both the symmetric and antisymmetric blocks of the full policy-FOC Jacobian, together with strict regime, curvature, feasibility and welfare-gap margins, establishes a common nonempty open primitive neighborhood in which this policy configuration remains a global SPNE. None of these steps establishes uniqueness of \emph{all} government equilibria.

\begin{proposition}[Canonical global-SPNE bridge]
\label{prop:global-spne}
At the canonical rational primitives
\[
(k_x,k_g,\mu,\nu,\kappa,d,e,\beta,\xi,m,\bar E)
=(4,18,0.9,1.2,0.8,2,1.1,1.7,0.1,2,0),
\]
the symmetric interior-duopoly policy candidate is a global policy Nash equilibrium for every $\theta\in[0,1]$ once the complete nonnegative Stage-2 and Stage-3 continuations are taken into account. Together with the unique continuation in every downstream subgame, this policy profile induces an SPNE of the full game. Moreover, the strict branch and welfare-gap inequalities, the continuous endpoint scaling of the exclusionary gap, and policy-FOC Jacobian regularity imply that the same local global-SPNE configuration survives on a nonempty open neighborhood of the canonical primitives.
\end{proposition}

A separate, finite rational certificate is required for the main fixed-rivalry comparison in Proposition~\ref{prop:fixed-theta-target}. At $\theta=31/40$ and $d=39/20,41/20$, Appendix~\ref{app:fixed-target-global} re-solves the symmetric government policies, checks their positive duopoly continuations and full government-FOC Jacobians, and independently maximizes all five full-policy continuation regions. At both points, rival-kink and rival-monopoly deviations are unreachable by nonnegative policies; the best own-kink/own-monopoly deviations meet at a common boundary and fall short by strictly positive rational welfare gaps. Strictness and regularity suffice for local comparative-static derivatives of the \emph{global} policy best response under small rival-policy changes. They do not certify every target weight between these values or uniqueness of every government equilibrium.

Proposition~\ref{prop:global-spne} is a qualification bridge, not a general primitive characterization of global equilibrium. The hostile-audit counterexample retained in the repository demonstrates that the interior Hessian and quartic conditions alone do not guarantee the proposition away from the certified region.